% Generated draft: requires mathematical and source review.
\documentclass[11pt]{article}
\usepackage[margin=25mm]{geometry}
\usepackage{fontspec}
\setmainfont{Latin Modern Roman}
\usepackage{amsmath,amssymb,mathtools,unicode-math}
\setmathfont{Latin Modern Math}
\usepackage{enumitem,needspace,xurl,hyperref}
\setlength{\parindent}{0pt}
\setlength{\parskip}{5pt}
\setlength{\emergencystretch}{3em}
\newcommand{\N}{\mathbb{N}}
\newcommand{\Z}{\mathbb{Z}}
\newcommand{\Q}{\mathbb{Q}}
\newcommand{\R}{\mathbb{R}}
\newcommand{\C}{\mathbb{C}}
\newcommand{\F}{\mathbb{F}}
\newcommand{\A}{\mathbb{A}}
\newcommand{\PP}{\mathbb{P}}
\newcommand{\E}{\mathbb{E}}
\newcommand{\eps}{\varepsilon}
\newcommand{\dd}{\,\mathrm{d}}
\begin{document}
\section*{UnsolvedMath ID: 30004712}
\textbf{Source problem code:} catalogue entry
\subsection{Definitions and mathematical statement}
Write \(\pi_U:M_U\to U\) for the given smooth projective fibration and \(i:U\hookrightarrow\mathbb{P}^n\) for the supplied inclusion. A compactification of \(M_U\) means a compact space \(M\) with an open immersion \(j:M_U\hookrightarrow M\) having dense image, in the geometric category intended by the source.

\textbf{Compactification question.} Characterize the given fibrations \(\pi_U\) for which there exist a smooth hyper-Kähler compactification \((M,j)\) and an everywhere-defined regular morphism \(\pi:M\to\mathbb{P}^n\) satisfying
\[
\pi\circ j=i\circ\pi_U.
\]
The requested characterization is a necessary-and-sufficient condition on the input fibration; the record does not supply such a condition.

\textbf{Unresolved conventions.} The source does not specify the ground field, the geometric structure or openness of \(U\), the allowed values of \(n\), or whether ``fibration'' requires surjectivity and connected fibres. The usual complex-geometric interpretation would take the ground field to be \(\C\), but this is an interpretation requiring confirmation. For a compact complex manifold, an irreducible holomorphic symplectic interpretation of ``hyper-Kähler'' requires that \(M\) be Kähler and simply connected, with
\[
H^0(M,\Omega_M^2)=\C\omega,
\qquad
\omega\text{ everywhere nondegenerate}.
\]
The record does not specify whether this irreducible convention or a broader hyper-Kähler convention is intended. It also does not explicitly require \(M\) to be projective. If \(U\) is open, the potentially stronger requirement \(j(M_U)=\pi^{-1}(U)\) is not explicitly stated and must be clarified before it is imposed. Thus the displayed extension condition records the supplied target, while these missing conventions prevent a unique fully specified formulation.
\subsection{Short English statement}
Determine exactly when a given smooth projective fibration over a subset of projective space admits a smooth hyper-Kähler compactification whose projection extends regularly to all of projective space. The intended geometric category, hyper-Kähler convention, and precise relation over \(U\) require clarification.
\subsection{Sources}
\begin{itemize}
\item Supplied local UnsolvedMath registry record, permanent ID \(30004712\), source problem code \texttt{catalogue entry}, entitled ``Hyperkähler Compactification of Projective Fibrations.'' The project identifies \texttt{data/unsolvedmath/problems.json} and the supplementary \texttt{catalog/problems.json} as local registries; their contents have not been independently inspected for this draft.
\item Supplied primary report citation: \emph{Classical Algebraic Geometry (2022)}, Oberwolfach Reports, DOI 10.4171/owr/2021/29, pp.\ 1519--1577. \url{https://doi.org/10.4171/owr/2021/29}. The background describes this collection as workshop reports documenting mathematical discussions and open problems. The record gives proposed year 2021 and supplies no named proposer.
\item Supplied literature reference, attributed to Saccà: \url{https://arxiv.org/abs/2411.06505}. No title or more specific bibliographic attribution is supplied.
\item Supplied literature reference, attributed to Saccà: \url{https://arxiv.org/abs/2002.01420}. No title or more specific bibliographic attribution is supplied.
\item Literature assessment dated exactly \texttt{2026-08-22}: the supplied curation labels the problem ``partially solved'' and states that compactification criteria and broad classes of torsors and group schemes are known, including smooth hyper-Kähler examples, while no necessary-and-sufficient criterion for the general question is established. This is supplied dated literature triage, not an independently verified assessment. The background expressly requests independent verification of time-sensitive status and cited claims.
\item The registry labels the statement \texttt{corrected\_verified}; its accompanying verification note describes only consistent mathematical delimiters for the projective-space inclusions and morphisms, without changing the question. This label does not establish mathematical verification of the formulation.
\end{itemize}
\end{document}
